\label{sec:resumo_plano}

O projeto, com vigência de 1º de setembro de 2025 a 31 de agosto de 2026, dava continuidade à etapa anterior. Nela, uma rede neural densa de \num{8195} parâmetros estimou o vértice de interação de depósitos de \qty{10}{\MeV} no terceiro módulo do detector distante (FD3) do experimento DUNE (\textit{Deep Underground Neutrino Experiment}) a partir de \num{60} descritores da luz de cintilação \cite{vicentin2024study}.

O objetivo era usar aprendizado de máquina para estimar a posição tridimensional de depósitos de energia da ordem de MeV a partir da luz registrada por um sistema de detecção de fótons do tipo APEX (\textit{Aluminum Profiles with Embedded X-ARAPUCA}), como medida do desempenho desse sistema. Previa-se medir a resolução espacial em cada coordenada e compará-la com a dos estimadores analíticos do estudo APEX \cite{marinho2025apex}. A reconstrução seria também estendida a outras energias depositadas e a outras geometrias de sensores, como propunha a conclusão da etapa anterior.

A metodologia partia da simulação usada naquela etapa, com as contagens de fótons de \qty{100}{\MeV} reescaladas para \qty{10}{\MeV} e detecção binomial com eficiência de \qty{2}{\percent}. O plano incluía a busca de hiperparâmetros, a avaliação em eventos reservados para teste pelo erro absoluto médio (MAE), pela raiz do erro quadrático médio (RMSE) e pelo coeficiente de determinação ($R^2$) em cada eixo, e um estudo em função da energia e da eficiência. Esperava\discretionary{-}{-}{-}se uma resolução por coordenada melhor que a da etapa anterior, cujo MAE de teste ia de \num{12.03} a \qty{18.18}{\cm} por eixo \cite{vicentin2024study}, e melhor que a dos estimadores analíticos, com um modelo leve o bastante para uso em aquisição em tempo real. Os resultados seriam apresentados no Encontro de Iniciação Científica do Instituto Tecnológico de Aeronáutica (ENCITA) de 2026 e neste relatório.
